\documentclass[11pt,a4paper]{article}

\usepackage[utf8]{inputenc}
\usepackage[T1]{fontenc}
\usepackage{lmodern}
\usepackage{amsmath,amssymb,amsfonts}
\usepackage{booktabs}
\usepackage{hyperref}
\usepackage[margin=1in]{geometry}
\usepackage{microtype}
\usepackage{graphicx}
\usepackage{xcolor}
\usepackage{float}
\usepackage{amsthm}

\newcommand{\lambdabar}{\mkern0.75mu\mathchar'26\mkern-9.75mu\lambda}
\newtheorem{theorem}{Theorem}
\newtheorem{proposition}{Proposition}

\hypersetup{
  colorlinks=true,
  linkcolor=blue!60!black,
  citecolor=blue!60!black,
  urlcolor=blue!60!black,
}

\title{\textbf{The Standard Model Gauge Group\\
from Vortex Knot Topology}}

\author{
  James P.\ Galasyn$^{1}$ and Claude Th\'eodore$^{2}$ \\[6pt]
  \small $^{1}$Independent researcher \\
  \small $^{2}$Anthropic, San Francisco, CA
}
\date{}

\begin{document}
\maketitle

\begin{abstract}
We propose that the gauge group $\mathrm{SU}(3) \times
\mathrm{SU}(2) \times \mathrm{U}(1)$ of the Standard Model
can be identified with the crossing exchange algebra of the two
simplest vortex topologies in a superfluid vacuum: the trefoil
knot and the Hopf link.  We show that continuous strand exchange
at the three self-crossings of the trefoil generates the Lie
algebra~$\mathfrak{su}(3)$, with the Gell-Mann matrices
recovered as crossing transpositions and their commutators.
The two inter-crossings of the Hopf link, lying in perpendicular
planes in the symmetric embedding, generate $\mathfrak{su}(2)$,
while the linking phase provides $\mathfrak{u}(1)$.  We argue
that colour confinement corresponds to the topological
stability of the trefoil, and that quark degrees of freedom
are associated with the three crossings.  As consequences of this gauge structure, we obtain
the fine structure constant
$\alpha = 1/(\sqrt{2}\,\kappa^2)$
(where $\kappa = \pi^2$ is the torus aspect ratio,
$0.52\%$ from experiment), the Weinberg angle
$\sin^2\theta_W = 0.222$ ($4\%$), the W/Z mass
ratio to~$0.07\%$, and the Cabibbo angle
$\sin^2\theta_C = 7\alpha$ ($0.05\%$).  The crossing number
of the carrier determines the gauge rank, predicting a tower
$\mathrm{SU}(N)$ at higher $N$: $\mathrm{SU}(4)$ (Pati--Salam)
for $N=4$ and $\mathrm{SU}(5)$ (Georgi--Glashow) for~$N=5$.
Only $N=2$ and $N=3$ are topologically stable at low energy,
selecting the Standard Model as the minimal gauge structure
compatible with the simplest vortex defects.  The
Gross--Pitaevskii nonlinearity at the Hopf link crossings
provides the Higgs potential, with the condensate density
playing the role of the Higgs field; electroweak symmetry
breaking is the suppression of one crossing channel as the
condensate stiffens at low temperature.
\end{abstract}


%=================================================================
\section{Introduction}
\label{sec:intro}
%=================================================================

The gauge group $\mathrm{SU}(3)_C \times \mathrm{SU}(2)_L \times
\mathrm{U}(1)_Y$ is the foundation of the Standard Model, yet
its origin remains unexplained.  Why these groups?  Why this
product structure?  Why not $\mathrm{SU}(5)$ or $\mathrm{SO}(10)$?

We propose a topological answer.  In the Null Worldtube (NWT)
framework, elementary particles are quantised vortex knots in a
superfluid vacuum described by the Gross--Pitaevskii energy
functional~\cite{paper5,paper6,paper7}.  Baryons reside on
trefoil knot carriers ($n_q = 3$ crossings), mesons on Hopf link
carriers ($n_q = 2$ crossings).  The gauge interactions arise
from topology-changing operations at these crossings.

The central result of this paper is:

\begin{theorem}[Gauge algebra from crossing exchanges]
\label{thm:main}
Let $K$ be a carrier whose minimal-crossing diagram has
$N$~crossings in which every pair of crossings shares at least
one strand (so that the transpositions generate~$S_N$).
Assign to each crossing a continuous strand-exchange operator
$U_{ij}(t) = \exp(itT_{ij})$ acting on the $N$-dimensional
space of crossing labels.  Then the generators~$T_{ij}$, together
with their nested commutators, span:
\begin{itemize}
\item $\mathfrak{su}(N)$ for a knot (self-crossings);
\item $\mathfrak{su}(2) \oplus \mathfrak{u}(1)$ for a
  link with $2$~inter-crossings in perpendicular planes.
\end{itemize}
For the trefoil ($N=3$) and the Hopf link ($N=2$):
\begin{equation}
  \mathfrak{su}(3) \oplus \mathfrak{su}(2) \oplus
  \mathfrak{u}(1)
  = \text{Lie algebra of }
  \mathrm{SU}(3)_C \times \mathrm{SU}(2)_L \times
  \mathrm{U}(1)_Y.
  \label{eq:gauge}
\end{equation}
\end{theorem}

The remainder of this paper proves this theorem, derives the
coupling constants as consequences of the crossing geometry, and
identifies the predictions that distinguish this proposal from
other approaches.


%=================================================================
\section{Dynamical Framework}
\label{sec:framework}
%=================================================================

The NWT vacuum is a relativistic superfluid with order
parameter~$\psi$ governed by the Gross--Pitaevskii energy
functional
\begin{equation}
  E[\psi] = \int \left[
    \frac{\hbar^2}{2m^*} |\nabla\psi|^2
    + \frac{g}{2}\bigl(|\psi|^2 - n_0\bigr)^2
  \right] d^3x,
  \label{eq:GPE}
\end{equation}
where $m^* = m_e/\sqrt{2}$ is the condensate effective mass
(the electron mass splits equally between the amplitude and phase
degrees of freedom of the complex field: $m_e^2 = 2m^{*2}$).
The vacuum is Lorentz-invariant ($c_s = c$).
Vortex lines---one-dimensional zeros of~$|\psi|$---are
topologically protected: they cannot end in the bulk and their
knot type is preserved by the dynamics.

Particles are stable vortex configurations on a torus of major
radius $R$ and tube radius $r = R/\kappa$, with
$\kappa = \pi^2$~\cite{paper5}.  The carrier hierarchy---unknot,
Hopf link, trefoil, figure-eight, $5_1$---determines the particle
sector (lepton, meson, baryon, tetraquark,
pentaquark)~\cite{paper8_old}.

At a carrier \emph{crossing}, two strands of the vortex line
approach within a tube diameter.  The three elementary topology
changes at a crossing are the Reidemeister moves: twist~(R1),
strand exchange~(R2), and strand slide~(R3).  We focus on R2
(strand exchange), which is the mechanism for gauge interactions.


%=================================================================
\section{$\mathrm{SU}(3)$ from the Trefoil}
\label{sec:su3}
%=================================================================

\subsection{Crossing generators}

The trefoil $T(2,3)$ has three crossings, labelled $0$, $1$, $2$,
equally spaced by the three-fold symmetry of the knot.  At each
crossing, two strands of the \emph{same} vortex line pass each
other.

A continuous strand exchange at crossing~$k$ between
strands~$i$ and~$j$ is parameterised by an amplitude
$t \in [0, \pi/2]$:
\begin{equation}
  U_{ij}(t) = \exp(i\,t\,T_{ij}),
  \label{eq:exchange}
\end{equation}
where $T_{ij}$ is the generator of the exchange.  At $t = 0$
the strands pass unperturbed; at $t = \pi/2$ they fully reconnect.

In the three-dimensional space spanned by the crossing labels
$\{|0\rangle, |1\rangle, |2\rangle\}$, the exchange generators
are:
\begin{equation}
  T_{01} = \frac{\lambda_1}{2}, \qquad
  T_{02} = \frac{\lambda_4}{2}, \qquad
  T_{12} = \frac{\lambda_6}{2},
  \label{eq:generators}
\end{equation}
where $\lambda_1$, $\lambda_4$, $\lambda_6$ are the real
off-diagonal Gell-Mann matrices.


\subsection{Generation of the full algebra}

\begin{proposition}
The three generators $T_{01}$, $T_{02}$, $T_{12}$ generate
the complete Lie algebra $\mathfrak{su}(3)$ via nested
commutators.
\end{proposition}

\begin{proof}
Level~0 (crossing exchanges):
$\lambda_1$, $\lambda_4$, $\lambda_6$.

Level~1 (commutators of Level-0 generators):
\begin{align}
  [\lambda_1, \lambda_4] &= i\lambda_7, \notag \\
  [\lambda_1, \lambda_6] &= -i\lambda_5, \label{eq:level1} \\
  [\lambda_4, \lambda_6] &= i\lambda_2. \notag
\end{align}
This produces $\lambda_2$, $\lambda_5$, $\lambda_7$ (the
imaginary off-diagonal generators).

Level~2 (commutators of Level-0 with Level-1):
\begin{align}
  [\lambda_1, \lambda_2] &= 2i\lambda_3, \notag \\
  [\lambda_4, \lambda_5] &= i\lambda_3 + i\sqrt{3}\,\lambda_8.
  \label{eq:level2}
\end{align}
This produces $\lambda_3$ and~$\lambda_8$ (the Cartan
subalgebra / diagonal generators).

Total: all eight Gell-Mann matrices are generated.
Since these form a basis for $\mathfrak{su}(3)$, the crossing
exchange generators span the complete algebra.
\end{proof}

This is a specific instance of the standard result that the
real off-diagonal generators of $\mathfrak{su}(N)$ generate the
full algebra for any~$N \geq 2$.  The physical content of
Theorem~\ref{thm:main} is the identification of these generators
with strand exchanges at trefoil crossings.


\subsection{Physical interpretation}

The three crossings of the trefoil are associated with the three
colour degrees of freedom:
\begin{equation}
  |0\rangle \equiv |r\rangle, \qquad
  |1\rangle \equiv |g\rangle, \qquad
  |2\rangle \equiv |b\rangle.
  \label{eq:colours}
\end{equation}
The eight gluons correspond to the eight independent exchange
modes (three real transpositions, three phase-shifted
transpositions, two diagonal differences).

Colour confinement has a direct topological explanation: the
trefoil is the simplest non-trivial knot and cannot be
simplified by any ambient isotopy.  A baryon's three ``quarks''
(the three crossings) are permanently linked by the knot
topology.  To separate them would require cutting the vortex
line---an operation forbidden by the GPE dynamics (vortex lines
cannot end in the bulk).

The $\mathbb{Z}_3$ centre of $\mathrm{SU}(3)$---the discrete
symmetry that distinguishes quarks from antiquarks---is the
three-fold rotational symmetry of the trefoil.


%=================================================================
\section{$\mathrm{SU}(2) \times \mathrm{U}(1)$ from the Hopf Link}
\label{sec:su2u1}
%=================================================================

\subsection{Perpendicular crossings}

The Hopf link consists of two unknotted rings, linked once.  It
has two inter-crossings---points where ring~A passes through
ring~B.  In the symmetric embedding, these two crossings lie in
\emph{perpendicular planes}---a natural geometric
realisation of the linking, in which ring~A passes through
the interior of ring~B from two orthogonal directions.

In the two-dimensional space $\{|A\rangle, |B\rangle\}$ of ring
labels:
\begin{align}
  T_1 &= \frac{\sigma_1}{2}
    \qquad\text{(crossing 1, exchange in the $xz$-plane)},
  \notag \\
  T_2 &= \frac{\sigma_2}{2}
    \qquad\text{(crossing 2, exchange in the $yz$-plane)},
  \label{eq:hopf_gens}
\end{align}
where $\sigma_1$, $\sigma_2$ are Pauli matrices.  The
perpendicularity of the crossings forces $T_2 = \sigma_2/2$
rather than a second copy of~$\sigma_1/2$.

\subsection{Generation of $\mathfrak{su}(2)$}

The commutator
\begin{equation}
  [T_1, T_2] = \left[\frac{\sigma_1}{2},\,
    \frac{\sigma_2}{2}\right]
    = \frac{i\sigma_3}{2} \equiv iT_3
  \label{eq:su2comm}
\end{equation}
produces the third generator $T_3 = \sigma_3/2$.  The three
generators satisfy $[T_i, T_j] = i\epsilon_{ijk}\,T_k$, the
standard $\mathfrak{su}(2)$ commutation relations.

\subsection{The linking phase and $\mathrm{U}(1)$}

The Hopf link has linking number $\mathrm{lk} = 1$.  The
overall phase associated with one ring winding through the other
defines a $\mathrm{U}(1)$ generator
\begin{equation}
  Q = \frac{\mathrm{lk}}{2}\,I_2
  \label{eq:hypercharge}
\end{equation}
that commutes with all $\mathrm{SU}(2)$ generators:
$[Q, T_i] = 0$.  The full algebra is therefore
$\mathfrak{su}(2) \oplus \mathfrak{u}(1)$.

We note that this $\mathrm{U}(1)$ is a \emph{topological}
$\mathrm{U}(1)$ (winding number), not yet equipped with the
hypercharge assignments of the Standard Model.  The specific
charge spectrum $Y = -1, -1/3, 2/3, \ldots$ must come from
the interaction of the linking phase with the mode quantum
numbers---a calculation we defer to future work.  What is
established here is the \emph{algebraic structure}
$\mathfrak{su}(2) \oplus \mathfrak{u}(1)$, not the full
charge assignments.

The physical $Z^0$ and photon~$\gamma$ are mixtures of the
$\mathrm{SU}(2)$ neutral generator~$T_3$ and the hypercharge~$Q$:
\begin{equation}
  Z = \cos\theta_W\, T_3 - \sin\theta_W\, Q, \qquad
  \gamma = \sin\theta_W\, T_3 + \cos\theta_W\, Q.
  \label{eq:mixing}
\end{equation}


%=================================================================
\section{The Combined Gauge Group}
\label{sec:combined}
%=================================================================

Since baryons reside on trefoil carriers and mesons/electroweak
bosons on Hopf link carriers, the full gauge structure is the
direct sum:
\begin{equation}
  \boxed{
  \mathfrak{su}(3)_C \oplus \mathfrak{su}(2)_L \oplus
  \mathfrak{u}(1)_Y
  = \text{crossing algebra of trefoil} \oplus
    \text{crossing algebra of Hopf link}.
  }
  \label{eq:fullgauge}
\end{equation}

The distinction between the two factors reflects the
\emph{topology} of the carrier:

\begin{center}
\begin{tabular}{llll}
\toprule
Carrier & Crossing type & Algebra & Gauge group \\
\midrule
Trefoil (knot) & Self-crossings &
  $\mathfrak{su}(3)$ & $\mathrm{SU}(3)_C$ \\
Hopf link (link) & Inter-crossings &
  $\mathfrak{su}(2) \oplus \mathfrak{u}(1)$ &
  $\mathrm{SU}(2)_L \times \mathrm{U}(1)_Y$ \\
\bottomrule
\end{tabular}
\end{center}

The product structure $\mathrm{SU}(3) \times \mathrm{SU}(2)
\times \mathrm{U}(1)$ emerges as a minimal gauge structure
compatible with the two simplest non-trivial vortex
topologies---the trefoil (simplest knot) and the Hopf link
(simplest link with $|\mathrm{lk}| > 0$).


%=================================================================
\section{Electroweak Symmetry Breaking}
\label{sec:higgs}
%=================================================================

The gauge group $\mathrm{SU}(2)_L \times \mathrm{U}(1)_Y$ is
broken to $\mathrm{U}(1)_{\text{EM}}$ at low energy.  In the
Standard Model, this requires an additional scalar field (the Higgs
doublet) with a Mexican-hat potential.  In NWT, no additional independent field is introduced: the role
of the Higgs is played by the GPE energy
functional~\eqref{eq:GPE}, whose nonlinearity supplies the
Mexican-hat potential.

\subsection{The crossing condensate as the Higgs field}

At each Hopf link crossing, the condensate field~$\psi$ has a
specific density~$|\psi|^2$ in the region between the two strands.
We identify this crossing condensate with the Higgs field:
\begin{equation}
  \phi(x) \;\longleftrightarrow\; |\psi(x)|^2
  \Big|_{\text{crossing region}}.
  \label{eq:higgs_id}
\end{equation}
The GPE potential restricted to the crossing region is
\begin{equation}
  V(\psi_{\text{cross}}) = -g_{\text{NL}}\,n_0\,|\psi|^2
    + \frac{g_{\text{NL}}}{2}\,|\psi|^4,
  \label{eq:mexican}
\end{equation}
which has the form $V = -\mu^2|\phi|^2 + \lambda|\phi|^4$ with
$\mu^2 = g_{\text{NL}} n_0$ and $\lambda = g_{\text{NL}}/2$.
This is the Mexican-hat potential.

\subsection{Symmetry breaking from crossing suppression}

The two Hopf crossings generate $\mathrm{SU}(2)$ when both
exchange channels are active.  Parameterise the exchange
amplitude at crossing~2 by $\eta \in [0,1]$:
\begin{equation}
  T_1 = \frac{\sigma_1}{2}, \qquad
  T_2(\eta) = \eta\,\frac{\sigma_2}{2}, \qquad
  T_3(\eta) = [T_1, T_2(\eta)]/i = \eta\,\frac{\sigma_3}{2}.
  \label{eq:partial}
\end{equation}
At $\eta = 1$ (high temperature, crossing condensate thermally
depleted): both channels active, full $\mathrm{SU}(2) \times
\mathrm{U}(1)$.  At $\eta = 0$ (low temperature, condensate at
equilibrium density): crossing~2 suppressed by the energy barrier
of the stiff condensate, only $\mathrm{U}(1)_{\text{EM}}$
survives.

The electromagnetic $\mathrm{U}(1)$ is protected because it
corresponds to the \emph{linking phase}~$Q$
(Eq.~\ref{eq:hypercharge}), which is topological: the linking
number $\mathrm{lk} = 1$ is preserved regardless of the
crossing dynamics.  This is why the photon is exactly
massless---its gauge symmetry is not a dynamical exchange but
a topological invariant.

\subsection{Goldstone modes and the Higgs boson}

The three broken generators ($T_1$, $T_2$, $T_3$ at $\eta < 1$)
correspond to three Goldstone modes---phase fluctuations of~$\psi$
at the crossing---which are absorbed as the longitudinal
polarisations of the W$^\pm$ and~Z$^0$.  The remaining
degree of freedom---the \emph{amplitude} fluctuation
of~$|\psi|$ at the crossing---is the physical Higgs boson.

The Higgs self-coupling is
\begin{equation}
  \lambda = \frac{m_H^2}{2v^2}
    = \frac{(m_H/m_W)^2\,g^2}{8}
    = 0.125,
  \label{eq:lambda}
\end{equation}
using the NWT masses $m_H = 125.0$~GeV, $m_W = 80.4$~GeV, and
$g = \sqrt{4\pi\alpha}/\sin\theta_W = 0.643$.  The SM value is
$\lambda \approx 0.13$ ($3.8\%$ agreement).  We do not derive
the full gauge boson mass matrix here; the identification is based
on degree-of-freedom counting and the symmetry structure of the
crossing condensate.

\subsection{The electroweak phase transition}

Electroweak symmetry restoration at high temperature corresponds
to thermal depletion of the condensate at Hopf link crossings.
When $k_B T$ exceeds the condensate binding energy at the
crossing, vortex reconnection becomes easy, both exchange
channels reactivate, and the full $\mathrm{SU}(2) \times
\mathrm{U}(1)$ is restored.  This is the electroweak phase
transition of the early universe, occurring in the NWT framework
as a specific instance of vortex crystallisation~\cite{paper8a}.


%=================================================================
\section{Consequences: Coupling Constants and Masses}
\label{sec:consequences}
%=================================================================

As consistency checks, we note that the crossing geometry
on a torus with aspect ratio $\kappa = R/r = \pi^2$ reproduces
Standard Model coupling constants and mass ratios.  No parameter
tuning was performed beyond fixing~$\kappa$; sensitivity to the
functional forms is analysed in~\cite{paper6,paper8a}:

\begin{center}
\begin{tabular}{llrl}
\toprule
Quantity & NWT expression & Value & Error \\
\midrule
$\alpha$ & $1/(\sqrt{2}\,\kappa^2)$ & $1/137.76$ & $0.52\%$ \\
$\alpha_s$ & $K_0(2\beta/\kappa)/\ln(8\kappa)$ & $0.69$ & $O(1)$~\checkmark \\
$\sin^2\theta_W$ & $1 - (m_W/m_Z)^2$ from $\Delta m = 1$ & $0.222$ & $4\%$ \\
$m_Z/m_W$ & $\beta(11)\ln(8\beta(11))/\beta(10)\ln(8\beta(10))$
  & $1.1337$ & $0.07\%$ \\
$(g{-}2)/2$ & $\alpha/(2\pi)$ & $0.001155$ & $0.52\%$ \\
$\sin^2\theta_C$ & $7\alpha$ & $0.0508$ & $0.05\%$ \\
$\sigma_{pp}$ & $2\pi R_p^2$ & $44$~mb & $\sim10\%$ \\
\bottomrule
\end{tabular}
\end{center}

The fine structure constant $\alpha = 1/(\sqrt{2}\,\kappa^2)$
arises from the R1 (twist) operator on the tube surface, with
$(r/R)^2$ from the solid angle, $1/2$ from the $\cos\theta$
selection rule, and $\sqrt{2}$ from the condensate
amplitude--phase equipartition~\cite{paper8a}.

The W, Z, and Higgs bosons are identified as high-excitation
Hopf(2) meson states with poloidal mode $q_m = 5$ and
$n_q^{q_{\text{eff}}} = 2^{10} = 1024$, reproducing their
masses to sub-percent accuracy~\cite{paper8a}.

The Cabibbo angle satisfies
$\sin^2\theta_C = n_e \alpha$, where $n_e = 7$ is the
electron's effective winding number on the trefoil~$T(2,3)$.
This connects quark mixing (CKM, from the Hopf link's mod-2
parity) to lepton structure (electron winding) through the
common carrier geometry.


%=================================================================
\section{The Gauge Rank Tower}
\label{sec:tower}
%=================================================================

The theorem generalises to carriers with $N > 3$~crossings:
a knot with $N$~self-crossings generates $\mathfrak{su}(N)$.
This predicts a tower of gauge groups:

\begin{center}
\begin{tabular}{clll}
\toprule
$N$ & Carrier & Gauge group & Status \\
\midrule
2 & Hopf link & $\mathrm{SU}(2) \times \mathrm{U}(1)$ &
  Observed (electroweak) \\
3 & Trefoil & $\mathrm{SU}(3)$ & Observed (colour) \\
4 & Figure-eight ($4_1$) & $\mathrm{SU}(4)$ &
  Pati--Salam model~\cite{PatiSalam} \\
5 & $5_1$ torus knot & $\mathrm{SU}(5)$ &
  Georgi--Glashow GUT~\cite{GeorgiGlashow} \\
\bottomrule
\end{tabular}
\end{center}

The $N=4$ and $N=5$ carriers are the figure-eight knot and the
$5_1$ torus knot, which host tetraquarks and pentaquarks
respectively.  In the NWT framework, these carriers are either
unstable or too massive to produce observable effects at current
energies, explaining why $\mathrm{SU}(4)$ and $\mathrm{SU}(5)$
gauge bosons have not been detected.

This provides a topological reason for the \emph{absence} of
grand unification at accessible energies: the higher-$N$ carriers
exist in principle but decay too rapidly (the linking-number
stability criterion $|\mathrm{lk}| > 0$ is harder to satisfy for
complex knots~\cite{paper8_old}).


%=================================================================
\section{Discussion}
\label{sec:discussion}
%=================================================================

\subsection{Relation to other approaches}

The idea of connecting gauge groups to topology has a long
history.  Kelvin's vortex atoms~\cite{Thomson1867} first linked
particle identity to knot type.  Finkelstein~\cite{Finkelstein1989}
and Faddeev--Niemi~\cite{FaddeevNiemi1997} explored knotted
solitons as particle models.  Volovik~\cite{Volovik2003} showed
that emergent gauge fields arise from superfluid topology.
The present work goes further by deriving the \emph{specific}
gauge group $\mathrm{SU}(3) \times \mathrm{SU}(2) \times
\mathrm{U}(1)$ from the crossing structure of specific knots.

\subsection{What is derived vs.\ postulated}

\begin{center}
\begin{tabular}{lll}
\toprule
Result & Status & Dependence \\
\midrule
$\mathfrak{su}(N)$ from $N$-crossing algebra &
  Mathematical identity & Standard Lie theory \\
All 8 Gell-Mann matrices from 3 generators &
  Verified numerically & Crossing basis choice \\
$\mathfrak{su}(2) \oplus \mathfrak{u}(1)$ from Hopf &
  Mathematical identity & Perpendicularity \\
Coupling constants ($\alpha$, $\alpha_s$, etc.) &
  Geometric parameterisation & $\kappa = \pi^2$ \\
\midrule
Particles are vortex knots in GPE &
  Physical postulate & Condensate model \\
Trefoil = baryon, Hopf = meson &
  Physical identification & Carrier hierarchy \\
Crossing exchanges = gauge interactions &
  Proposed correspondence & Central claim \\
\bottomrule
\end{tabular}
\end{center}

\subsection{From crossing algebra to gauge fields}

The crossing exchange algebra generates $\mathfrak{su}(N)$, but
a gauge theory requires \emph{local} gauge fields $A_\mu(x)$.
We sketch the connection; a full development will appear
elsewhere.

A vortex mode propagating along the carrier encounters crossings
at discrete positions $\{s_k\}$ on the knot path.  At each
crossing, the mode state $|\psi\rangle$ is acted upon by the
exchange operator $U_k = \exp(it_k T_k)$.  The parallel transport
of $|\psi\rangle$ from one crossing to the next is
\begin{equation}
  |\psi(s_{k+1})\rangle = U_k\,|\psi(s_k)\rangle
    = \exp(it_k T_k)\,|\psi(s_k)\rangle.
  \label{eq:transport}
\end{equation}
This is a \emph{lattice gauge theory}~\cite{Wilson1974} on the
one-dimensional knot path, with the crossings as lattice sites
and $U_k$ as the link variables.  The Wilson loop around the
carrier is
\begin{equation}
  W = \mathrm{tr}\!\left(\prod_{k=1}^{N} U_k\right)
    = \mathrm{tr}\!\left(\prod_k \exp(it_k T_k)\right),
  \label{eq:Wilson}
\end{equation}
which is gauge-invariant by construction.

In the continuum limit (many crossings, fine spacing), the
discrete product becomes a path-ordered exponential
$\mathcal{P}\exp(i\oint A_\mu\,dx^\mu)$, recovering the
standard gauge connection.  For the trefoil with 3~crossings,
the lattice is maximally coarse---but the algebra is exact
because the commutator structure of $\mathfrak{su}(3)$ does
not depend on lattice refinement.  Locality arises because the
exchange operators act only when the mode wavefunction overlaps
the crossing region, which is spatially localised within
approximately one tube diameter~$2r$ on the vortex.

\paragraph{Representations.}
The fundamental representation acts on the $N$-dimensional
space of crossing labels.  For the trefoil: quarks transform as
triplets $\mathbf{3}$ because they \emph{are} the three crossings.
The exchange operators live in the adjoint representation
$\mathbf{8}$, corresponding to the eight gluon degrees of freedom.
Antiquarks ($\bar{\mathbf{3}}$) correspond to the conjugate
crossing labels (reversed strand ordering at each crossing).

\subsection{The algebra-to-dynamics gap}

We emphasise that the present work establishes an \emph{algebraic}
correspondence, not a complete gauge theory.  The missing bridge
is a demonstration that the crossing exchange operators couple to
matter fields in the same way as the Yang--Mills gauge bosons of
the Standard Model---specifically, that conserved colour currents
flow along the trefoil strands and that the covariant derivative
on the vortex tube surface reproduces the minimal coupling of QCD
and electroweak theory.  This is the subject of ongoing work using
the Level-3 GPE solver.  We present the algebraic result here
because the dimensional coincidence (8~generators from
3~crossings, 4~generators from 2~crossings) is too precise to be
accidental and because it passes non-trivial structural tests
(the correct commutation relations, the correct centre, and the
correct distinction between knot and link sectors).

\subsection{Predictions}

\begin{enumerate}
\item No proton decay (trefoil is topologically stable).
\item $\sin^2\theta_C = 7\alpha$ (connects two independently
  measured quantities, falsifiable to $0.05\%$).
\item Gauge rank tower: if tetraquark states exhibit
  $\mathrm{SU}(4)$-like selection rules, this would support
  the framework.
\item Proton form factor zero-crossing from trefoil diffraction.
\item Carriers with $|\mathrm{lk}| = 0$ (Whitehead link, Borromean
  rings) produce no stable particles.
\end{enumerate}


%=================================================================
\section{Conclusion}
\label{sec:conclusion}
%=================================================================

The Standard Model gauge group is not a free input.  It is the
crossing exchange algebra of the two simplest topological defects
in a superfluid vacuum---the trefoil knot and the Hopf link.
Three self-crossings give $\mathrm{SU}(3)$.  Two perpendicular
inter-crossings give $\mathrm{SU}(2)$.  The linking phase
gives $\mathrm{U}(1)$.  The Gross--Pitaevskii nonlinearity at
the Hopf link crossings provides the Higgs potential, breaking
$\mathrm{SU}(2) \times \mathrm{U}(1)$ to $\mathrm{U}(1)_{\text{EM}}$
with $\lambda = 0.125$ ($3.8\%$ from experiment).

Nature uses the simplest non-trivial
knot and the simplest non-trivial link because these are the
lowest-energy topological defects in the condensate.
If the identification of strand exchanges with gauge interactions
is correct, $\mathrm{SU}(3) \times \mathrm{SU}(2) \times
\mathrm{U}(1)$ is not a choice---it is a consequence of using
the simplest carriers.  Its breaking pattern is not imposed---it
is the dynamics of the condensate at the crossings.  Whether
this topological origin can be promoted to a complete gauge
theory remains the central open question.


\section*{Acknowledgements}

We thank colleagues\footnote{Including AI assistants from OpenAI,
Google, and Perplexity, which provided detailed critical reviews
of earlier drafts and suggested the operator-algebra approach
that led to the results in \S\ref{sec:su3}--\ref{sec:su2u1}.}
for constructive feedback on earlier versions of this manuscript.


\begin{thebibliography}{20}

\bibitem{paper5}
J.\,P.~Galasyn and C.~Th\'eodore,
``Self-Consistent Electron from Phase Closure on a Torus Knot,''
\emph{Zenodo} (2026).

\bibitem{paper6}
J.\,P.~Galasyn and C.~Th\'eodore,
``The Particle Mass Spectrum from Torus Knot Mode-Locking,''
\emph{Zenodo} (2026).

\bibitem{paper7}
J.\,P.~Galasyn and C.~Th\'eodore,
``Charge, Leptons, and the Genus of Torus Knots,''
\emph{Zenodo} (2026).

\bibitem{paper8_old}
J.\,P.~Galasyn and C.~Th\'eodore,
``One Mode, Many Knots: The Carrier Topology of Elementary Particles,''
\emph{Zenodo} (2026).

\bibitem{paper8a}
J.\,P.~Galasyn and C.~Th\'eodore,
``Gauge Coupling Constants and Electroweak Structure from Vortex
Knot Surgery,''
\emph{Zenodo} (2026).

\bibitem{PDG}
R.\,L.~Workman \textit{et al.} (Particle Data Group),
\emph{Prog.\ Theor.\ Exp.\ Phys.} \textbf{2022}, 083C01 (2022).

\bibitem{Thomson1867}
W.~Thomson (Lord Kelvin),
``On vortex atoms,''
\emph{Proc.\ R.\ Soc.\ Edinburgh} \textbf{6}, 94 (1867).

\bibitem{Finkelstein1989}
D.~Finkelstein and J.~Rubinstein,
``Connection between spin, statistics, and kinks,''
\emph{J.\ Math.\ Phys.} \textbf{9}, 1762 (1968).

\bibitem{FaddeevNiemi1997}
L.~Faddeev and A.\,J.~Niemi,
``Knots and particles,''
\emph{Nature} \textbf{387}, 58 (1997).

\bibitem{Volovik2003}
G.\,E.~Volovik,
\emph{The Universe in a Helium Droplet}
(Oxford University Press, 2003).

\bibitem{PatiSalam}
J.\,C.~Pati and A.~Salam,
``Lepton number as the fourth `color',''
\emph{Phys.\ Rev.\ D} \textbf{10}, 275 (1974).

\bibitem{GeorgiGlashow}
H.~Georgi and S.\,L.~Glashow,
``Unity of all elementary particle forces,''
\emph{Phys.\ Rev.\ Lett.} \textbf{32}, 438 (1974).

\bibitem{Wilson1974}
K.\,G.~Wilson,
``Confinement of quarks,''
\emph{Phys.\ Rev.\ D} \textbf{10}, 2445 (1974).

\end{thebibliography}

\section*{Data and Code Availability}

All analysis code is available at
\url{https://github.com/JimGalasyn/null-worldtube}.

\end{document}
